% Options for packages loaded elsewhere
\PassOptionsToPackage{unicode}{hyperref}
\PassOptionsToPackage{hyphens}{url}
\documentclass[
]{article}
\usepackage{xcolor}
\usepackage[margin=1in]{geometry}
\usepackage{amsmath,amssymb}
\setcounter{secnumdepth}{-\maxdimen} % remove section numbering
\usepackage{iftex}
\ifPDFTeX
  \usepackage[T1]{fontenc}
  \usepackage[utf8]{inputenc}
  \usepackage{textcomp} % provide euro and other symbols
\else % if luatex or xetex
  \usepackage{unicode-math} % this also loads fontspec
  \defaultfontfeatures{Scale=MatchLowercase}
  \defaultfontfeatures[\rmfamily]{Ligatures=TeX,Scale=1}
\fi
\usepackage{lmodern}
\ifPDFTeX\else
  % xetex/luatex font selection
\fi
% Use upquote if available, for straight quotes in verbatim environments
\IfFileExists{upquote.sty}{\usepackage{upquote}}{}
\IfFileExists{microtype.sty}{% use microtype if available
  \usepackage[]{microtype}
  \UseMicrotypeSet[protrusion]{basicmath} % disable protrusion for tt fonts
}{}
\makeatletter
\@ifundefined{KOMAClassName}{% if non-KOMA class
  \IfFileExists{parskip.sty}{%
    \usepackage{parskip}
  }{% else
    \setlength{\parindent}{0pt}
    \setlength{\parskip}{6pt plus 2pt minus 1pt}}
}{% if KOMA class
  \KOMAoptions{parskip=half}}
\makeatother
\usepackage{graphicx}
\makeatletter
\newsavebox\pandoc@box
\newcommand*\pandocbounded[1]{% scales image to fit in text height/width
  \sbox\pandoc@box{#1}%
  \Gscale@div\@tempa{\textheight}{\dimexpr\ht\pandoc@box+\dp\pandoc@box\relax}%
  \Gscale@div\@tempb{\linewidth}{\wd\pandoc@box}%
  \ifdim\@tempb\p@<\@tempa\p@\let\@tempa\@tempb\fi% select the smaller of both
  \ifdim\@tempa\p@<\p@\scalebox{\@tempa}{\usebox\pandoc@box}%
  \else\usebox{\pandoc@box}%
  \fi%
}
% Set default figure placement to htbp
\def\fps@figure{htbp}
\makeatother
% definitions for citeproc citations
\NewDocumentCommand\citeproctext{}{}
\NewDocumentCommand\citeproc{mm}{%
  \begingroup\def\citeproctext{#2}\cite{#1}\endgroup}
\makeatletter
 % allow citations to break across lines
 \let\@cite@ofmt\@firstofone
 % avoid brackets around text for \cite:
 \def\@biblabel#1{}
 \def\@cite#1#2{{#1\if@tempswa , #2\fi}}
\makeatother
\newlength{\cslhangindent}
\setlength{\cslhangindent}{1.5em}
\newlength{\csllabelwidth}
\setlength{\csllabelwidth}{3em}
\newenvironment{CSLReferences}[2] % #1 hanging-indent, #2 entry-spacing
 {\begin{list}{}{%
  \setlength{\itemindent}{0pt}
  \setlength{\leftmargin}{0pt}
  \setlength{\parsep}{0pt}
  % turn on hanging indent if param 1 is 1
  \ifodd #1
   \setlength{\leftmargin}{\cslhangindent}
   \setlength{\itemindent}{-1\cslhangindent}
  \fi
  % set entry spacing
  \setlength{\itemsep}{#2\baselineskip}}}
 {\end{list}}
\usepackage{calc}
\newcommand{\CSLBlock}[1]{\hfill\break\parbox[t]{\linewidth}{\strut\ignorespaces#1\strut}}
\newcommand{\CSLLeftMargin}[1]{\parbox[t]{\csllabelwidth}{\strut#1\strut}}
\newcommand{\CSLRightInline}[1]{\parbox[t]{\linewidth - \csllabelwidth}{\strut#1\strut}}
\newcommand{\CSLIndent}[1]{\hspace{\cslhangindent}#1}
\setlength{\emergencystretch}{3em} % prevent overfull lines
\providecommand{\tightlist}{%
  \setlength{\itemsep}{0pt}\setlength{\parskip}{0pt}}
\usepackage{bookmark}
\IfFileExists{xurl.sty}{\usepackage{xurl}}{} % add URL line breaks if available
\urlstyle{same}
\hypersetup{
  pdftitle={Testing the Jelly Ocean Hypothesis in the Gulf of Maine},
  pdfauthor={Nicholas R. Record, Benjamin Tupper},
  hidelinks,
  pdfcreator={LaTeX via pandoc}}

\title{Testing the Jelly Ocean Hypothesis in the Gulf of Maine}
\author{Nicholas R. Record, Benjamin Tupper}
\date{}

\begin{document}
\maketitle

\subsection{Introduction}\label{introduction}

There has often been a ``crustaceocentric'' bias within zooplankton
ecology in the ocean {[}1{]}. The study of gelatinous species,
collectively termed gelata, has been underemphasized as compared to
crustaceans. In the early 2000s, gelata had a moment in the spotlight,
following conspicuous blooms that damaged fisheries {[}2{]}, aquaculture
(ref), energy production (Seshadri et al.~2000), and recreation (ref).
In addition to motivating new research on gelata, these events raised
the question of whether there is a larger pattern of change emerging.

The Jelly Ocean Hypothesis posits that gelatinous zooplankton
populations are increasing worldwide as a result of changing ocean
conditions, primarily due to anthropogenic activities. A series of
recent studies have argued that, due a multiplicity of causes, abundance
of gelatinous species is growing. Proposed causes are primarily
human-driven, including warming (ref), eutrophication, translocation
{[}2{]}, overfishing, bottom trawling (Qui 2014), pollution,
aquaculture, and acidification {[}3{]}. There is also evidence that
shifts toward dominance of gelata are difficult to reverse because of
dynamics like histeresis {[}4{]}.

Testing the Jelly Ocean Hypothesis requires long time series
(\textgreater20 years) of gelata abundances in order to compare to
historical levels. The shortage of such datasets has left the Jelly
Ocean Hypothesis largely unsupported {[}5{]}. In lieu of time series,
studies have used qualitative data such as news reports to argue that
gelata are increasing globally {[}6{]}. However, it is difficult to
distinguish these perceived increases from multi-annual or multi-decadal
fluctuations. A recent meta-analysis of time series showed that most
perceived increases reflect a \textasciitilde20 year cyclicity, with a
worldwide upswing in the 1990s, coinciding with the increase in concern
over gelata {[}7{]}, although some notable long-term increases in gelata
were not included in the analysis {[}8{]}. The Jelly Ocean Hypothesis
remains very much in debate {[}9{]}.

In some ways, the Jelly Ocean Hypothesis is loosely posed. Gelata are
taxonomically very diverse, encompassing multiple phyla. This diversity
confounds attempts to quantify increasing trends. The list of species
classified under the gelatinous umbrella varies, and includes taxa with
drastic and important differences in their life histories, trophic
positions, and ecological roles. For example, pelagic ctenophores are
primarily copepod predators, where as pelagic tunicates graze some of
the smallest microbes. Cnidarians and ctenophores are the two most
common taxa to group together in analysis, and yet they are
phylogenetically very distant (Moroz et al.~2014). Perhaps the only
characteristic common to the many gelatinous phyla is a body with a high
water content, alternatively expressed as a low carbon:volume ratio.
There is no a priori reason to expect that phyla with such striking
differences and perhaps only a superficial similarity should be grouped
together with respect to their temporal trends. Some have suggested that
the Jelly Ocean Hypothesis is ill-posed---i.e.~which species are
increasing, where, and for what reason (Gibbons \& Richardson 2013)?

Despite their major differences, there is some evidence that a diverse
collection of gelata can vary together (Condon et al.~2013). This
suggests that certain conditions can occur that favor the gelatinous
body form in general, over other strategies that utilize a high
carbon:volume ratio (e.g.~crustaceans). This reframes the Jelly Ocean
Hypothesis in trait-based terms: conditions in the ocean ecosystem can
shift to benefit the gelatinous body form. This would be inclusive of
gelata in general, despite differences in phylogeny, trophic position,
or other life history traits.

We revisited the Jelly Ocean Hypothesis using a 50-year spatial survey
on the northwest North Atlantic Shelf. We addressed two main questions:
1) what are the temporal changes in gelata presence; and 2) what are the
underlying conditions associated with changes. We examined three
taxonomic groups well-represented in the dataset--siphonophores, salps,
and coelenterates--to explore the idea that certain conditions can favor
the gelatinous body plan in general, rather a particulars.

\subsection{Methods}\label{methods}

\subsubsection{Data}\label{data}

The Northeast Fisheries Science Center Ecosystem Monitoring (ECOMON)
program collects seasonal zooplankton measurements as part of its
mission to monitor and assess the Northeast Continental Shelf ecosystem.
Sampling includes oblique tows using a 61 cm bongo net fitted with a 333
µm mesh at stratified locations across the shelf (Figure 1). Tows are a
minimum of five minutes in duration and sample from the surface to
within 5 m of the seabed or to a maximum depth of 200 m. A mechanical
flowmeter fitted in the mouth of each net is used to quantify volume
sampled. Gelatinous zooplankton are identified into broad taxonomic
groups, including siphonophores, ctenophores, salps, and coelenterates
(i.e.~scyphozoans). The broad grouping, while missing \ldots{} has
allowed for internal consistency over the duration of the survey. The
analysis conducted here uses presence or absence of each taxonomic
group.

\subsubsection{Time series analysis}\label{time-series-analysis}

\subsubsection{Spatial modeling}\label{spatial-modeling}

\subsection{Results}\label{results}

\subsection{Discussion}\label{discussion}

\subsection*{References}\label{references}
\addcontentsline{toc}{subsection}{References}

\phantomsection\label{refs}
\begin{CSLReferences}{0}{1}
\bibitem[\citeproctext]{ref-haddock2004golden}
1. Haddock SH: \textbf{A golden age of gelata: Past and future research
on planktonic ctenophores and cnidarians}. \emph{Hydrobiologia} 2004,
\textbf{530}:549--556.

\bibitem[\citeproctext]{ref-kideys2005impacts}
2. Kideys AE, Roohi A, Bagheri S, Finenko G, Kamburska L:
\textbf{Impacts of invasive ctenophores on the fisheries of the black
sea and caspian sea}. \emph{Oceanography} 2005, \textbf{18}:76--85.

\bibitem[\citeproctext]{ref-richardson2008jellyfish}
3. Richardson AJ, Gibbons MJ: \textbf{Are jellyfish increasing in
response to ocean acidification?} \emph{Limnology and Oceanography}
2008, \textbf{53}:2040--2045.

\bibitem[\citeproctext]{ref-richardson2009jellyfish}
4. Richardson AJ, Bakun A, Hays GC, Gibbons MJ: \textbf{The jellyfish
joyride: Causes, consequences and management responses to a more
gelatinous future}. \emph{Trends in ecology \& evolution} 2009,
\textbf{24}:312--322.

\bibitem[\citeproctext]{ref-condon2012questioning}
5. Condon RH, Graham WM, Duarte CM, Pitt KA, Lucas CH, Haddock SH,
Sutherland KR, Robinson KL, Dawson MN, Decker MB, et al.:
\textbf{Questioning the rise of gelatinous zooplankton in the world's
oceans}. \emph{Bioscience} 2012, \textbf{62}:160--169.

\bibitem[\citeproctext]{ref-brotz2012increasing}
6. Brotz L, Cheung WW, Kleisner K, Pakhomov E, Pauly D:
\textbf{Increasing jellyfish populations: Trends in large marine
ecosystems}. In \emph{Jellyfish blooms IV: Interactions with humans and
fisheries}. Springer; 2012:3--20.

\bibitem[\citeproctext]{ref-condon2013recurrent}
7. Condon RH, Duarte CM, Pitt KA, Robinson KL, Lucas CH, Sutherland KR,
Mianzan HW, Bogeberg M, Purcell JE, Decker MB, et al.: \textbf{Recurrent
jellyfish blooms are a consequence of global oscillations}.
\emph{Proceedings of the National Academy of Sciences} 2013,
\textbf{110}:1000--1005.

\bibitem[\citeproctext]{ref-atkinson2004long}
8. Atkinson A, Siegel V, Pakhomov E, Rothery P: \textbf{Long-term
decline in krill stock and increase in salps within the southern ocean}.
\emph{Nature} 2004, \textbf{432}:100--103.

\bibitem[\citeproctext]{ref-sanz2016flawed}
9. Sanz-Martı́n M, Pitt KA, Condon RH, Lucas CH, Novaes de Santana C,
Duarte CM: \textbf{Flawed citation practices facilitate the
unsubstantiated perception of a global trend toward increased jellyfish
blooms}. \emph{Global Ecology and Biogeography} 2016,
\textbf{25}:1039--1049.

\end{CSLReferences}

\end{document}
